\section{Introduction}

Advances in large language models have catalysed a shift toward agentic systems capable of multi-step planning, tool use, and autonomous action~\cite{AbouAli2025}. In regulated deployment contexts---public administration, healthcare, legal services---this shift introduces a distinctive accountability challenge: as agents act with greater independence, the chain of responsibility between a decision and a human decision-maker grows longer and more opaque. Consequential actions become harder to attribute, audit trails harder to maintain, and errors more difficult to reverse before they propagate.

Two design decisions jointly determine how a deployment handles these challenges. \emph{Autonomy} describes how much human involvement a workflow requires, from a human who commands every step to one who monitors an autonomously operating system. \emph{Agency} describes what the system can perceive and affect, from a model that reasons over supplied text to one that writes to authoritative records. Prior work treats the two as independent design dimensions~\cite{Feng2025Levels}, yet in regulated deployment they cannot be navigated in isolation: a highly capable agent without appropriate oversight may breach accountability norms; one too tightly constrained offers negligible practical benefit~\cite{Mitchell2025}.

We address this gap by introducing a two-dimensional design space---five autonomy levels and five agency levels---and six architectural tactics for navigating it. The analysis is grounded in an automated funding program matching system, demonstrating how the tactics adjust a deployment's configuration under realistic compliance constraints. The result is a shared vocabulary for reasoning about where oversight sits, what consequences agent actions can have, and how reversibility is preserved---considerations that are central to accountable deployment in regulated contexts.
